\chapter{अवकलन और उत्तलता}\label{ch:g12:deriv}

\omterm{def:g12:deriv:derivative}{अवकलज} किसी \omterm{def:g11:func:function}{फलन} के परिवर्तन की तात्क्षणिक दर नापता है; वह उसके \omterm{def:g10:functions:graph}{आलेख} की \omterm{def:g12:deriv:tangent}{स्पर्श
रेखा} की \omterm{def:g10:reffunc:affine}{प्रवणता} है। यह अध्याय अवकलन के नियमों को दोहराता और आगे बढ़ाता है,
$f'$ के चिह्न को $f$ के परिवर्तन से जोड़ता है, और \omterm{def:g12:deriv:convex}{उत्तलता} प्रस्तुत करता है,
जिस पर द्वितीय \omterm{def:g12:deriv:derivative}{अवकलज} का राज है।

\section{अवकलज}

\begin{definition}[किसी बिंदु पर अवकलज]\label{def:g12:deriv:derivative}
मान लीजिए $f$ किसी \omterm{def:g10:numbers:interval}{अंतराल} $I$ पर परिभाषित है और $a \in I$ है। \omterm{def:g11:func:function}{फलन} $f$ \emph{$a$
पर अवकलनीय}\index{अवकलज}\index{अवकलनीय} है यदि अंतर-भागफल
\[
\frac{f(a+h) - f(a)}{h}
\]
की $h \to 0$ पर कोई परिमित सीमा हो। इस सीमा को \emph{$a$ पर $f$ का अवकलज} कहते
हैं और $f'(a)$ लिखते हैं। यदि $f$ $I$ के हर बिंदु पर अवकलनीय हो, तो \omterm{def:g11:func:function}{फलन} $f' \colon x \mapsto f'(x)$
$f$ का \emph{अवकलज} कहलाता है।
\end{definition}

\begin{definition}[स्पर्श रेखा]\label{def:g12:deriv:tangent}
यदि $f$ $a$ पर \omterm{def:g12:deriv:derivative}{अवकलनीय} हो, तो बिंदु $(a, f(a))$ पर $f$ के वक्र की \emph{स्पर्श
रेखा}\index{स्पर्श रेखा} इस \omterm{def:g10:algebra:equation}{समीकरण} वाली रेखा है
\[
y = f(a) + f'(a)(x - a).
\]
\end{definition}

\begin{omfigure}
\begin{tikzpicture}
\begin{axis}[omaxis, clip=false,
  width=9.6cm, height=6.8cm,
  xmin=-0.4, xmax=3.4, ymin=-0.4, ymax=4.3,
  xtick={1,2.5}, xticklabels={$a$,$a+h$}, ytick=\empty]
  \addplot[omProp!40, thick, domain=0.2:2.7] {1 + 1.75*(x-1)};
  \addplot[omProp!70, thick, domain=0.07:2.95] {1 + 1.5*(x-1)};
  \addplot[omProp, thick, domain=-0.12:3.3] {1 + 1.25*(x-1)};
  \addplot[omThm, thick, domain=-0.4:3.3] {x};
  \addplot[omDef, thick, domain=-0.4:2.72] {x^2/2 + 0.5};
  \draw[black, dashed, thin] (axis cs:1,1) -- (axis cs:2.5,1)
    node[midway, below, font=\small] {$h$};
  \draw[black, dashed, thin] (axis cs:2.5,1) -- (axis cs:2.5,3.625)
    node[midway, right, font=\small] {$f(a+h)-f(a)$};
  \fill (axis cs:2.5,3.625) circle (1.5pt);
  \fill (axis cs:2,2.5) circle (1.5pt);
  \fill (axis cs:1.5,1.625) circle (1.5pt);
  \fill (axis cs:1,1) circle (1.5pt)
    node[above left, font=\small] {$A$};
  \node[omThm, font=\small, anchor=west] at (axis cs:3.38,3.3) {स्पर्श रेखा};
\end{axis}
\end{tikzpicture}

{\small जैसे-जैसे $h$ घटता है, $A = (a, f(a))$ और $(a+h, f(a+h))$ से होकर जाने वाली \omterm{def:g11:deriv:rate}{छेदक रेखा}
(नारंगी) $A$ पर की \omterm{def:g12:deriv:tangent}{स्पर्श रेखा} (लाल) की ओर घूमती जाती है: उसकी \omterm{def:g10:reffunc:affine}{प्रवणता} $\frac{f(a+h)-f(a)}{h}$
$f'(a)$ की ओर जाती है।}
\end{omfigure}

\begin{proposition}\label{prop:g12:deriv:diffcont}
$a$ पर \omterm{def:g12:deriv:derivative}{अवकलनीय} \omterm{def:g11:func:function}{फलन} $a$ पर \omterm{def:g12:limcont:continuity}{संतत} होता है। इसका विलोम असत्य है।
\end{proposition}

\begin{proof}
छोटे $h \neq 0$ के लिए $f(a+h) - f(a) = h \cdot \frac{f(a+h)-f(a)}{h}$। $h \to 0$ पर दाहिना पक्ष $0 \times f'(a) = 0$ की ओर जाता है, इसलिए $f(a+h) \to f(a)$।
विलोम के लिए: $x \mapsto \abs{x}$ $0$ पर \omterm{def:g12:limcont:continuity}{संतत} है, पर $0$ पर उसका अंतर-भागफल $h$ के चिह्न
के अनुसार $\frac{\abs h}{h} = \pm 1$ के बराबर होता है, और उसकी कोई सीमा नहीं है।
\end{proof}

\subsection{अवकलन के नियम}

\begin{proposition}[संक्रियाएँ]\label{prop:g12:deriv:operations}
मान लीजिए $u, v$ $I$ पर \omterm{def:g12:deriv:derivative}{अवकलनीय} हैं और $\lambda \in \R$ है। तब $u + v$, $\lambda u$, $uv$ $I$ पर
\omterm{def:g12:deriv:derivative}{अवकलनीय} हैं, और जहाँ $v \neq 0$ है वहाँ $\frac{u}{v}$ भी, तथा
\[
(u+v)' = u' + v', \quad (\lambda u)' = \lambda u', \quad
(uv)' = u'v + uv', \quad
\left(\frac{u}{v}\right)' = \frac{u'v - uv'}{v^2}.
\]
\end{proposition}

\begin{proof}[गुणनफल के नियम की उपपत्ति]
$a$ पर $uv$ का अंतर-भागफल इस रूप में लिखिए:
\[
\frac{u(a+h)v(a+h) - u(a)v(a)}{h}
= \frac{u(a+h) - u(a)}{h}\,v(a+h) + u(a)\,\frac{v(a+h) - v(a)}{h}.
\]
$h \to 0$ पर \omterm{def:g12:limcont:continuity}{सांतत्य} से $v(a+h) \to v(a)$ (\cref{prop:g12:deriv:diffcont}), इसलिए दाहिना पक्ष $u'(a)v(a) + u(a)v'(a)$ की ओर जाता है। शेष
नियम इसी तरह सिद्ध होते हैं।
\end{proof}

\begin{theorem}[शृंखला नियम]\label{thm:g12:deriv:chain}
\index{शृंखला नियम}
मान लीजिए $u$ $I$ पर \omterm{def:g12:deriv:derivative}{अवकलनीय} है और उसके मान $J$ में हैं, तथा $g$ $J$ पर
\omterm{def:g12:deriv:derivative}{अवकलनीय} है। तब $g \circ u$ $I$ पर \omterm{def:g12:deriv:derivative}{अवकलनीय} है और
\[
(g \circ u)' = (g' \circ u) \cdot u' .
\]
विशेष रूप से, \omterm{def:g12:deriv:derivative}{अवकलनीय} $u$ के लिए:
\[
(u^n)' = n\,u'\,u^{n-1} \ (n \in \Z), \qquad
\left(\sqrt{u}\right)' = \frac{u'}{2\sqrt{u}} \ (u > 0), \qquad
\left(\eu^{u}\right)' = u'\,\eu^{u}.
\]
\end{theorem}

\begin{proof}[उपपत्ति की रूपरेखा]
जब छोटे $h$ के लिए $u(a+h) \neq u(a)$ हो, तब लिखिए
\[
\frac{g(u(a+h)) - g(u(a))}{h}
= \frac{g(u(a+h)) - g(u(a))}{u(a+h) - u(a)} \cdot \frac{u(a+h) - u(a)}{h}.
\]
$h \to 0$ पर $u(a+h) \to u(a)$, इसलिए पहला गुणनखंड $g'(u(a))$ की ओर और दूसरा $u'(a)$ की ओर जाता है।
(पूरी उपपत्ति को उस स्थिति को भी सँभालना पड़ता है जहाँ मनमाने छोटे $h$ के लिए
$u(a+h) = u(a)$ हो; इस तकनीकी बिंदु का उपचार स्नातक स्तर पर होता है।)
\end{proof}

\begin{example}
अपने-अपने स्वाभाविक \omterm{def:g11:func:function}{प्रांतों} पर सही, सामान्य \omterm{def:g12:deriv:derivative}{अवकलज}:
\[
(x^n)' = nx^{n-1}, \quad
\left(\frac{1}{x}\right)' = -\frac{1}{x^2}, \quad
(\sqrt{x})' = \frac{1}{2\sqrt{x}}, \quad
(\cos x)' = -\sin x, \quad
(\sin x)' = \cos x .
\]
पहला गुणनफल के नियम से आगमन द्वारा सिद्ध होता है, और अंतिम दो \cref{ch:g12:trigo} में।
\end{example}

\section{परिवर्तन और चरम मान}

\begin{theorem}[अवकलज का चिह्न और परिवर्तन]\label{thm:g12:deriv:variations}
मान लीजिए $f$ किसी \omterm{def:g10:numbers:interval}{अंतराल} $I$ पर \omterm{def:g12:deriv:derivative}{अवकलनीय} है।
\begin{enumerate}
  \item $f$ $I$ पर \omterm{def:g12:seq:monotonic}{वर्धमान} है यदि और केवल यदि $I$ पर $f' \geq 0$ हो।
  \item $f$ $I$ पर \omterm{def:g10:functions:variations}{ह्रासमान} है यदि और केवल यदि $I$ पर $f' \leq 0$ हो।
  \item यदि $I$ पर $f' > 0$ हो, सिवाय सीमित बिंदुओं के जहाँ वह लुप्त होता है,
        तो $f$ $I$ पर निरंतर \omterm{def:g12:seq:monotonic}{वर्धमान} है।
\end{enumerate}
\end{theorem}

\begin{proof}[आंशिक उपपत्ति]
यदि $f$ \omterm{def:g12:seq:monotonic}{वर्धमान} हो, तो हर अंतर-भागफल $\frac{f(a+h)-f(a)}{h}$ $\geq 0$ होता है, और सीमा लेने पर
$f'(a) \geq 0$ मिलता है। विलोम निष्कर्ष माध्यमान प्रमेय पर टिके हैं, जो स्नातक स्तर
पर सिद्ध होती है; उन्हें \emph{इस स्तर पर स्वीकार कर लिया जाता है}।
\end{proof}

\begin{definition}[स्थानीय चरम मान]\label{def:g12:deriv:extremum}
$f$ का $a$ पर \emph{स्थानीय अधिकतम}\index{चरम मान} है यदि $a$ के पास
के सभी $x$ के लिए $f(x) \leq f(a)$ हो; स्थानीय \omterm{def:g10:functions:extrema}{न्यूनतम} सममित रूप से परिभाषित होता है।
\end{definition}

\begin{proposition}[प्रथम-कोटि की शर्त]\label{prop:g12:deriv:critical}
यदि $f$ किसी \emph{खुले} \omterm{def:g10:numbers:interval}{अंतराल} $I$ पर \omterm{def:g12:deriv:derivative}{अवकलनीय} हो और $a \in I$ पर उसका कोई
स्थानीय \omterm{def:g12:deriv:extremum}{चरम मान} हो, तो $f'(a) = 0$। इसका विलोम असत्य है ($a = 0$ पर $f(x) = x^3$)।
\end{proposition}

\begin{proof}
मान लीजिए $a$ एक \omterm{def:g12:deriv:extremum}{स्थानीय अधिकतम} है। छोटे $h > 0$ के लिए $\frac{f(a+h)-f(a)}{h} \leq 0$, इसलिए $h \to 0^+$
लेने पर $f'(a) \leq 0$ मिलता है; और $h < 0$ के लिए भागफल $\geq 0$ होता है, जिससे $f'(a) \geq 0$
मिलता है। अतः $f'(a) = 0$।
\end{proof}

\begin{omfigure}
\begin{tikzpicture}
\begin{axis}[omaxis,
  width=8.4cm, height=5.6cm,
  xmin=-2.5, xmax=2.5, ymin=-4.2, ymax=4.2,
  xtick={-1,1}, ytick={-2,2}]
  \addplot[omProp, thick, domain=-1.8:-0.2] {2};
  \addplot[omProp, thick, domain=0.2:1.8] {-2};
  \addplot[omDef, thick, domain=-2.15:2.15] {x^3 - 3*x};
  \fill (axis cs:-1,2) circle (1.5pt);
  \fill (axis cs:1,-2) circle (1.5pt);
\end{axis}
\end{tikzpicture}

{\small किसी खुले \omterm{def:g10:numbers:interval}{अंतराल} के भीतर स्थानीय \omterm{def:g12:deriv:extremum}{चरम मान} पर \omterm{def:g12:deriv:tangent}{स्पर्श रेखा} (नारंगी)
क्षैतिज होती है: $f'(a) = 0$। यहाँ $f(x) = x^3 - 3x$, जिसमें $-1$ पर \omterm{def:g12:deriv:extremum}{स्थानीय अधिकतम} और $1$
पर स्थानीय \omterm{def:g10:functions:extrema}{न्यूनतम} है।}
\end{omfigure}

\begin{method}[परिवर्तन-सारणी]\label{met:g12:deriv:table}
किसी \omterm{def:g11:func:function}{फलन} $f$ का अध्ययन करने के लिए: उसका \omterm{def:g11:func:function}{प्रांत} और सीमाओं पर मान निकालिए;
$f'$ निकालकर उसका \omterm{def:g10:algebra:expand}{गुणनखंडन} कीजिए; हर उप-अंतराल पर $f'$ का चिह्न तय कीजिए;
सब कुछ \omterm{def:g10:functions:table}{परिवर्तन-सारणी} में दर्ज कीजिए और \omterm{def:g12:deriv:extremum}{चरम मान} अंकित कीजिए; और एकैकी आच्छादन
प्रमेय (\cref{thm:g12:limcont:bijection}) से $f(x) = k$ के हलों की संख्या निकालिए।
\end{method}

\section{उत्तलता}

\begin{definition}[उत्तल फलन]\label{def:g12:deriv:convex}
किसी \omterm{def:g10:numbers:interval}{अंतराल} $I$ पर परिभाषित \omterm{def:g11:func:function}{फलन} $f$ $I$ पर
\emph{उत्तल}\index{उत्तलता} है यदि उसके \omterm{def:g10:functions:graph}{आलेख} की हर जीवा \omterm{def:g10:functions:graph}{आलेख} के ऊपर पड़े:
अर्थात् सभी $a, b \in I$ और $t \in \intcc{0}{1}$ के लिए
\[
f\bigl((1-t)a + t b\bigr) \leq (1-t) f(a) + t f(b).
\]
वह \emph{अवतल}\index{अवतल} है यदि उल्टी असमिका सही हो ($-f$ उत्तल)।
\end{definition}

\begin{omfigure}
\begin{tikzpicture}
\begin{axis}[omaxis,
  width=8.8cm, height=6cm,
  xmin=-1.9, xmax=2.5, ymin=-0.8, ymax=4.3,
  xtick={-1,1.5}, xticklabels={$a$,$b$}, ytick=\empty]
  \addplot[omThm, thick, domain=-1.25:1.75] {0.5*x + 1.5};
  \addplot[omProp, thick, domain=-0.45:2.3] {x - 0.25};
  \addplot[omDef, thick, domain=-1.6:2.0] {x^2};
  \draw[black, dashed, thin] (axis cs:0.25,0.0625) -- (axis cs:0.25,1.625);
  \fill (axis cs:-1,1) circle (1.5pt);
  \fill (axis cs:1.5,2.25) circle (1.5pt);
  \fill (axis cs:0.5,0.25) circle (1.5pt);
  \node[omThm, font=\small, above left] at (axis cs:1.7,2.4) {जीवा};
  \node[omProp, font=\small, below right] at (axis cs:1.75,1.5) {स्पर्श रेखा};
\end{axis}
\end{tikzpicture}

{\small एक \omterm{def:g12:deriv:convex}{उत्तल} \omterm{def:g11:func:function}{फलन}: हर जीवा (लाल) \omterm{def:g10:functions:graph}{आलेख} के ऊपर पड़ती है, और \omterm{def:g10:functions:graph}{आलेख} अपनी हर
\omterm{def:g12:deriv:tangent}{स्पर्श रेखा} (नारंगी) के ऊपर। बिंदुदार रेखाखंड $f\bigl((1-t)a+tb\bigr)$ और $(1-t)f(a)+tf(b)$ के बीच का अंतर
दिखाता है।}
\end{omfigure}

\begin{theorem}[उत्तलता और अवकलज]\label{thm:g12:deriv:convexity}
मान लीजिए $f$ किसी \omterm{def:g10:numbers:interval}{अंतराल} $I$ पर \omterm{def:g12:deriv:derivative}{अवकलनीय} है। ये कथन तुल्य हैं:
\begin{enumerate}
  \item $f$ $I$ पर \omterm{def:g12:deriv:convex}{उत्तल} है;
  \item $f'$ $I$ पर \omterm{def:g12:seq:monotonic}{वर्धमान} है;
  \item $f$ का \omterm{def:g10:functions:graph}{आलेख} अपनी हर \omterm{def:g12:deriv:tangent}{स्पर्श रेखा} के ऊपर पड़ता है।
\end{enumerate}
यदि $f$ दो बार \omterm{def:g12:deriv:derivative}{अवकलनीय} हो, तो ये $I$ पर $f'' \geq 0$ के भी तुल्य हैं।
\end{theorem}

\begin{proof}[(2) $\Rightarrow$ (3) की उपपत्ति]
$a \in I$ स्थिर रखिए और $g(x) = f(x) - f(a) - f'(a)(x-a)$ रखिए, जो $a$ पर वक्र और \omterm{def:g12:deriv:tangent}{स्पर्श रेखा} के बीच का
अंतर है। तब $g$ \omterm{def:g12:deriv:derivative}{अवकलनीय} है और $g'(x) = f'(x) - f'(a)$। यदि $f'$ \omterm{def:g12:seq:monotonic}{वर्धमान} हो, तो $I \cap \intoc{-\infty}{a}$ पर
$g' \leq 0$ और $I \cap \intco{a}{+\infty}$ पर $g' \geq 0$: $g$ $a$ से पहले घटता है और बाद में बढ़ता है,
इसलिए $g$ अपना \omterm{def:g10:functions:extrema}{न्यूनतम} $g(a) = 0$ पर पाता है, अतः $g \geq 0$। शेष निष्कर्ष इस स्तर
पर स्वीकार कर लिए जाते हैं; $f'' \geq 0$ के साथ तुल्यता $f'$ पर \cref{thm:g12:deriv:variations} लगाने से
निकलती है।
\end{proof}

\begin{definition}[नति-परिवर्तन बिंदु]\label{def:g12:deriv:inflection}
जिस बिंदु पर $f$ का वक्र अपनी \omterm{def:g12:deriv:tangent}{स्पर्श रेखा} को पार कर जाता है वह
\emph{नति-परिवर्तन बिंदु}\index{नति-परिवर्तन बिंदु} है। दो बार \omterm{def:g12:deriv:derivative}{अवकलनीय} $f$
के लिए नति-परिवर्तन बिंदु वही हैं जहाँ $f''$ चिह्न बदलता है।
\end{definition}

\begin{example}\label{ex:g12:deriv:cube}
$f(x) = x^3$ के लिए $f''(x) = 6x$: $f$ $\intoc{-\infty}{0}$ पर \omterm{def:g12:deriv:convex}{अवतल} और $\intco{0}{+\infty}$ पर \omterm{def:g12:deriv:convex}{उत्तल} है, और मूल बिंदु
पर उसका \omterm{def:g12:deriv:inflection}{नति-परिवर्तन बिंदु} है, जहाँ वक्र अपनी \omterm{def:g12:deriv:tangent}{स्पर्श रेखा} ($x$-अक्ष) को
पार कर जाता है।
\end{example}

\begin{omfigure}
\begin{tikzpicture}
\begin{axis}[omaxis,
  width=7.6cm, height=5.6cm,
  xmin=-1.7, xmax=1.7, ymin=-2.9, ymax=2.9,
  xtick={-1,1}, ytick={-1,1}]
  \addplot[omProp, thick, domain=-1.2:1.2] {0};
  \addplot[omDef, thick, domain=-1.4:1.4] {x^3};
  \fill (axis cs:0,0) circle (1.5pt);
  \node[font=\small, anchor=east] at (axis cs:-0.45,-1.7) {अवतल};
  \node[font=\small, anchor=west] at (axis cs:0.45,1.7) {उत्तल};
\end{axis}
\end{tikzpicture}

{\small मूल बिंदु पर $x \mapsto x^3$ का \omterm{def:g12:deriv:inflection}{नति-परिवर्तन बिंदु}: वक्र अपनी \omterm{def:g12:deriv:tangent}{स्पर्श रेखा}
(नारंगी) को पार कर जाता है, बाईं ओर \omterm{def:g12:deriv:convex}{अवतल} और दाईं ओर \omterm{def:g12:deriv:convex}{उत्तल}।}
\end{omfigure}

\begin{example}[उत्तलता से असमिकाएँ]\label{ex:g12:deriv:ineq}
चरघातांकी \omterm{def:g11:func:function}{फलन} $\R$ पर \omterm{def:g12:deriv:convex}{उत्तल} है (उसका द्वितीय \omterm{def:g12:deriv:derivative}{अवकलज} वही स्वयं है, जो धनात्मक
है)। $0$ पर उसकी \omterm{def:g12:deriv:tangent}{स्पर्श रेखा} $y = 1 + x$ है, इसलिए
\[
\eu^x \geq 1 + x \quad x \in \R \text{ के लिए} .
\]
ऐसी \emph{स्पर्श-रेखा असमिकाएँ} \omterm{def:g12:deriv:convex}{उत्तलता} की मानक उपज हैं।
\end{example}

\begin{method}[उत्तलता का उपयोग]\label{met:g12:deriv:convexity}
\begin{itemize}
  \item कोई असमिका $f(x) \geq ax + b$ सिद्ध करने के लिए: उस रेखा को किसी \omterm{def:g12:deriv:convex}{उत्तल} $f$ की
        \omterm{def:g12:deriv:tangent}{स्पर्श रेखा} के रूप में पहचानिए और \cref{thm:g12:deriv:convexity}(3) का हवाला दीजिए।
  \item \omterm{def:g12:deriv:inflection}{नति-परिवर्तन बिंदु} ढूँढ़ने के लिए: $f''(x) = 0$ हल कीजिए \emph{और} जाँचिए
        कि वहाँ $f''$ चिह्न बदलता है।
  \item \omterm{def:g10:functions:table}{परिवर्तन-सारणी} में \omterm{def:g12:deriv:convex}{उत्तलता} वक्र के रेखाचित्र को और पैना कर देती है
        (वह किस ओर मुड़ता है)।
\end{itemize}
\end{method}

\section{अभ्यास}

\begin{exercise}[$\star$]\label{exo:g12:deriv:1}
निम्नलिखित \omterm{def:g10:functions:function}{फलनों} का उनके \omterm{def:g11:func:function}{प्रांतों} पर अवकलन कीजिए:
\[
f(x) = x^3 - 5x + 2, \quad
g(x) = \frac{x}{x^2+1}, \quad
h(x) = (2x+1)^5, \quad
k(x) = \sqrt{x^2 + 1}.
\]
\end{exercise}

\begin{exercise}[$\star$]\label{exo:g12:deriv:2}
\omterm{def:g10:coordgeom:system}{भुज} $a = 2$ वाले बिंदु पर $f(x) = x^2 - 3x + 1$ के वक्र की \omterm{def:g12:deriv:tangent}{स्पर्श रेखा} का \omterm{def:g10:algebra:equation}{समीकरण} बताइए, और इस
\omterm{def:g12:deriv:tangent}{स्पर्श रेखा} के सापेक्ष वक्र की स्थिति ज्ञात कीजिए।
\end{exercise}

\begin{exercise}[$\star$]\label{exo:g12:deriv:3}
\omterm{def:g12:deriv:derivative}{अवकलज} की परिभाषा का उपयोग करके निकालिए:
\[
\lim_{x \to 0} \frac{(1+x)^{100} - 1}{x}
\qquad\text{और}\qquad
\lim_{h \to 0} \frac{\sqrt{4+h} - 2}{h}.
\]
\end{exercise}

\begin{exercise}[$\star\star$]\label{exo:g12:deriv:4}
\omterm{def:g11:func:function}{फलन} $f(x) = \dfrac{x^2 + 3}{x - 1}$ का उसके \omterm{def:g11:func:function}{प्रांत} पर अध्ययन कीजिए: सीमाएँ, \omterm{def:g12:deriv:derivative}{अवकलज}, \omterm{def:g10:functions:table}{परिवर्तन-सारणी},
स्थानीय \omterm{def:g12:deriv:extremum}{चरम मान}।
\end{exercise}

\begin{exercise}[$\star\star$]\label{exo:g12:deriv:5}
भुजा $30\,$ cm की एक वर्गाकार गत्ते की चादर के कोनों से भुजा $x$ के बराबर
वर्ग काटकर और किनारे मोड़कर बिना ढक्कन का डिब्बा बनाया जाता है। डिब्बे के आयतन
को \omterm{def:g10:functions:extrema}{अधिकतम} करने वाला $x$ ज्ञात कीजिए।
\end{exercise}

\begin{exercise}[$\star\star$]\label{exo:g12:deriv:6}
मान लीजिए $f(x) = x^4 - 4x^3 + 10$ है।
\begin{enumerate}
  \item $f''$ निकालिए और $f$ के उत्तलता-अंतराल तथा \omterm{def:g12:deriv:inflection}{नति-परिवर्तन बिंदु}
        ज्ञात कीजिए।
  \item दिखाइए कि \omterm{def:g10:coordgeom:system}{भुज} $2$ वाले \omterm{def:g12:deriv:inflection}{नति-परिवर्तन बिंदु} पर \omterm{def:g12:deriv:tangent}{स्पर्श रेखा} वक्र को
        वहीं पार करती है।
\end{enumerate}
\end{exercise}

\begin{exercise}[$\star\star$]\label{exo:g12:deriv:7}
स्पर्श-रेखा असमिका का उपयोग करके दिखाइए कि सभी $x > 0$ के लिए
\[
\sqrt{x} \leq \frac{x + 1}{2},
\]
और बताइए कि समता कब होती है। (संकेत: वर्गमूल \omterm{def:g12:deriv:convex}{अवतल} है।)
\end{exercise}

\begin{exercise}[$\star\star\star$]\label{exo:g12:deriv:8}
मान लीजिए $f$ $\R$ पर \omterm{def:g12:deriv:convex}{उत्तल} और \omterm{def:g12:deriv:derivative}{अवकलनीय} है, और मान लीजिए $f'$ किसी बिंदु
$a$ पर लुप्त होता है। दिखाइए कि $f(a)$ $\R$ पर $f$ का \emph{वैश्विक}
\omterm{def:g10:functions:extrema}{न्यूनतम} है।
\end{exercise}

\begin{exercise}[$\star\star\star$]\label{exo:g12:deriv:9}
दिखाइए कि नियत परिमाप $p$ वाले सभी आयतों में वर्ग का क्षेत्रफल सबसे बड़ा
होता है। फिर दिखाइए कि नियत क्षेत्रफल $A$ वाले सभी आयतों में वर्ग का परिमाप
सबसे छोटा होता है, और समझाइए कि दोनों कथन आपस में कैसे जुड़े हैं।
\end{exercise}

\section{समस्या: जीवनरक्षक की गणना}

\begin{problem}[{सप्ताहांत समस्या --- सबसे तेज़ बचाव-मार्ग प्रकाश के स्नेल
नियम को मानता है, उत्तलता हर इष्टतम को प्रमाणित कर देती है, और आदर्श पेय-डिब्बा
ठीक उतना ही ऊँचा होता है जितना चौड़ा}]
\label{pb:g12:deriv:1}
एक जीवनरक्षक किसी तैराक को मुसीबत में देखती है --- तट के तिरछे नीचे, पानी के
भीतर। दौड़ना तैरने से तेज़ है: इसलिए सीधी रेखा सबसे तेज़ रास्ता \emph{नहीं}
है। समय को \omterm{def:g10:functions:extrema}{न्यूनतम} करने वाला रास्ता तट पर मुड़ जाता है, और उसके मुड़ने का नियम
ठीक वही है जिससे प्रकाश पानी में घुसते समय अपवर्तित होता है: प्रकृति भी अवकलन
करती है। यह समस्या शृंखला नियम (\cref{thm:g12:deriv:chain}) का अभ्यास कराती है, बचाव करती है,
\omterm{def:g12:deriv:convex}{उत्तलता} (\cref{thm:g12:deriv:convexity}) की असमिकाएँ काटती है, और एक टिन का डिब्बा बनाती है।

\medskip
\noindent\textbf{भाग I --- शृंखला नियम में प्रवाह।}
\begin{enumerate}
  \item अवकलन कीजिए: $\left(3x^2 + 1\right)^5$; \ $\sqrt{x^2 + 9}$; \ $\dfrac{1}{x^2 + 1}$।
  \item $x = 4$ पर $y = \sqrt{x^2 + 9}$ की \omterm{def:g12:deriv:tangent}{स्पर्श रेखा} बताइए।
  \item $\R$ पर $f(x) = \dfrac{x}{x^2 + 1}$ की \omterm{def:g10:functions:table}{परिवर्तन-सारणी} बनाइए (\omterm{def:g12:deriv:extremum}{चरम मान} सहित)।
  \item $g(x) = x^3 - 3x^2 + 4$ के लिए: \omterm{def:g10:functions:table}{परिवर्तन-सारणी} \emph{और} उत्तलता-सारणी ($g''$),
        \omterm{def:g12:deriv:inflection}{नति-परिवर्तन बिंदु} सहित (\cref{def:g12:deriv:inflection})।
  \item $g$ की उसके \omterm{def:g12:deriv:inflection}{नति-परिवर्तन बिंदु} पर \omterm{def:g12:deriv:tangent}{स्पर्श रेखा} निकालिए, और दिखाइए
        कि $g(x) - (\text{स्पर्श रेखा}) = (x - 1)^3$: चिह्न का यह बदलना इस बारे में क्या कहता है कि
        \omterm{def:g12:deriv:inflection}{नति-परिवर्तन बिंदु} पर \omterm{def:g12:deriv:tangent}{स्पर्श रेखा} वक्र से कैसे मिलती है?
\end{enumerate}

\medskip
\noindent\textbf{भाग II --- बचाव।}
तट $x$-अक्ष के अनुदिश है; जीवनरक्षक रेत पर $A(0, 30)$ पर खड़ी है, और तैराक $B(40, -20)$
पर प्रतीक्षा कर रहा है (मीटर में)। वह रेत पर $5$ m/s से दौड़ती है, $2$
m/s से तैरती है, और अपनी पसंद के किसी बिंदु $(x, 0)$ पर पानी में उतरती है।
\begin{enumerate}[resume]
  \item प्रतिरूप: कुल बचाव-समय $T(x)$ व्यक्त कीजिए, और बताइए कि केवल $x \in \intcc{0}{40}$
        ही विचार के योग्य क्यों है।
  \item $T$ का अवकलन कीजिए (शृंखला नियम काम पर --- दो बार)।
  \item मान लीजिए $\theta_1$ दौड़ वाले चरण का तट के \emph{लंब} से कोण है, और
        $\theta_2$ तैरने वाले चरण का। दिखाइए कि $\sin\theta_1 = \frac{x}{\sqrt{x^2 + 900}}$, और यह कि शर्त $T'(x) = 0$ का
        अर्थ है
        \[
        \frac{\sin\theta_1}{5} = \frac{\sin\theta_2}{2}
        \]
        --- यानी \emph{स्नेल का नियम}, जिसमें हवा और पानी में प्रकाश की
        चालों के स्थान पर ये दोनों चालें हैं।
  \item सीधा रेखाखंड $[AB]$ तट को $x = 24$ पर काटता है। $T(24)$, $T(40)$ (तट के
        साथ-साथ दौड़कर फिर सीधे तैरना) और $T(0)$ निकालिए: क्या ज्यामितीय सीधी
        रेखा सबसे तेज़ है? क्या दोनों चरम रणनीतियों में से कोई है?
  \item $T'$ पर \omterm{thm:g12:limcont:ivt}{द्विभाजन} से इष्टतम प्रवेश-बिंदु ढूँढ़िए (वह $\approx 33.7$ m है):
        $x^*$ मीटर तक और रिकॉर्ड समय $T(x^*)$ सेकंड के दसवें भाग तक बताइए।
        इष्टतम मोड़ सीधी रेखा की तुलना में कितना बचा लेता है?
  \item क्रांतिक बिंदु \emph{वैश्विक} \omterm{def:g10:functions:extrema}{न्यूनतम} क्यों है? ($T$ का हर पद
        \omterm{def:g12:deriv:convex}{उत्तल} है --- यह स्वीकार कर लीजिए कि \omterm{def:g12:deriv:convex}{उत्तल} \omterm{def:g10:functions:function}{फलनों} का योग \omterm{def:g12:deriv:convex}{उत्तल} होता
        है --- और \cref{exo:g12:deriv:8} लगाइए।)
  \item फ़र्मा का सिद्धांत कहता है कि प्रकाश सदा \omterm{def:g10:functions:extrema}{न्यूनतम} \emph{समय} वाला
        रास्ता लेता है। इससे निकालिए कि प्रकाश की किरण पानी में घुसते समय
        (जहाँ प्रकाश धीमा है) ठीक सतह पर क्यों मुड़ती है, और रोज़मर्रा का वह
        प्रेक्षण बताइए जिसकी व्याख्या इससे होती है --- पानी के गिलास में
        रखी पेंसिल।
\end{enumerate}

\medskip
\noindent\textbf{भाग III --- \omterm{def:g12:deriv:convex}{उत्तलता} की फ़सल।}
\begin{enumerate}[resume]
  \item दिखाइए कि $x \mapsto x^2$ \omterm{def:g12:deriv:convex}{उत्तल} है, और \omterm{prop:g10:coordgeom:midpoint}{मध्यबिंदु} असमिका $\left(\frac{a + b}{2}\right)^2 \leq
        \frac{a^2 + b^2}{2}$ दो बार सिद्ध
        कीजिए: एक बार \omterm{def:g12:deriv:convex}{उत्तलता} से, और एक बार $(a - b)^2 \geq 0$ का प्रसार करके।
  \item इससे निकालिए कि \emph{वर्ग-माध्य} $\sqrt{\frac{a^2 + b^2}{2}}$ \omterm{def:g12:seq:arith-geom}{समांतर} \omterm{def:g11:stat:mean}{माध्य} पर हावी रहता
        है, और इस शृंखला की पूरी कड़ी जोड़िए --- हरात्मक $\leq$ \omterm{def:g12:seq:arith-geom}{गुणोत्तर} $\leq$
        \omterm{def:g12:seq:arith-geom}{समांतर} $\leq$ वर्ग --- और चारों की जाँच $a = 2$, $b = 8$ पर कीजिए।
  \item स्पर्श-रेखा-नीचे-वक्र: $x \mapsto \frac1x$ $\intoo{0}{+\infty}$ पर \omterm{def:g12:deriv:convex}{उत्तल} है; $1$ पर उसकी \omterm{def:g12:deriv:tangent}{स्पर्श
        रेखा} लिखिए और सभी $x > 0$ के लिए असमिका $\frac1x \geq 2 - x$ निकालिए। समता कहाँ
        होती है?
  \item जंगल में नति-परिवर्तन: किसी महामारी में संचयी रोगियों की गिनती
        S-आकार के वक्र का अनुसरण करती है। उसके \omterm{def:g12:deriv:inflection}{नति-परिवर्तन बिंदु} पर क्या
        होता है, और महामारी-विज्ञानी उसी तारीख़ की प्रतीक्षा क्यों करते हैं?
        उसे दोनों ओर द्वितीय \omterm{def:g12:deriv:derivative}{अवकलज} के चिह्न से जोड़िए।
\end{enumerate}

\medskip
\noindent\textbf{भाग IV --- आदर्श डिब्बा।}
किसी बेलनाकार डिब्बे में कम से कम धातु के साथ $V = 330$ cm$^3$ समाना चाहिए:
पृष्ठ $S = 2\pi r^2 + 2\pi r h$, प्रतिबंध $\pi r^2 h = V$।
\begin{enumerate}[resume]
  \item $S(r) = 2\pi r^2 + \dfrac{2V}{r}$ व्यक्त कीजिए, अवकलन कीजिए, और $V = 330$ के लिए इष्टतम त्रिज्या
        तथा ऊँचाई निकालिए (मिलीमीटर की यथार्थता तक)।
  \item सुंदर सामान्य नियम सिद्ध कीजिए: इष्टतम पर $h = 2r$ --- आदर्श डिब्बा ठीक
        उतना ही ऊँचा होता है जितना चौड़ा।
  \item $S''(r) > 0$ जाँचिए और (फिर \cref{exo:g12:deriv:8} से) निष्कर्ष निकालिए कि इष्टतम वैश्विक
        है। असली पेय-डिब्बे साफ़ दिखता है कि चौड़ाई से अधिक ऊँचे होते हैं:
        कोई एक अ-गणितीय कारण बताइए।
  \item समापन --- इष्टतमीकरण की प्रक्रिया-शृंखला: प्रतिरूप बनाइए, \omterm{def:g12:deriv:derivative}{अवकलज} से
        पूछिए, क्रांतिक बिंदुओं के लिए हल कीजिए, \omterm{def:g12:deriv:convex}{उत्तलता} या सारणी से
        प्रमाणित कीजिए, और व्याख्या कीजिए। इस सूची को जीवनरक्षक, डिब्बे और
        कक्षा 11 के सबसे मज़बूत शहतीर (\cref{pb:g11:deriv:1}) पर चलाइए --- और बताइए कि
        फ़र्मा का सिद्धांत यह प्रक्रिया चलाने वाले और किसके बारे में कहता है।
\end{enumerate}
\end{problem}
